\section{Introduction}
\label{sec:intro}
% Sources: Results (audited), workshop paper text (unsure.txt: "Evaluation was restricted to
% in-distribution test data"; claimed larger gains for the lower-performing model; "evaluate only the
% aggregated consensus prediction rather than uncertainty maps"; no MC-Dropout/ensembles; no SegFormer),
% Wang 2019 abstract, Moshkov 2020 conclusions, Shanmugam 2021 abstract, Reinke 2024 abstract.

Test-time augmentation (TTA) runs a trained model on several transformed copies of an input, maps the outputs back and averages them. It needs no retraining, fits any segmentation pipeline, and tends to raise mean segmentation accuracy \cite{wang2019aleatoric,moshkov2020tta}; the spread of its predictions is also used as an uncertainty estimate \cite{wang2019aleatoric}. Because it is cheap and rarely hurts the average, it is easy to switch on without asking what else it changes; nnU-Net, for example, applies mirroring at test time by default \cite{isensee2021nnunet}.

Averages can hide the errors that matter in the clinic. An empty prediction adds a Dice of zero to the mean and has no boundary distance at all, so a method can raise mean Dice while losing lesions. In classification, TTA can change many correct predictions into incorrect ones even when accuracy improves overall \cite{shanmugam2021better}, and validation metrics in image analysis are often chosen inadequately \cite{maierhein2024metrics,reinke2024pitfalls}. For a method that will be switched on in deployment, the questions are what it does to missed lesions and false alarms on data from other centres, and whether its uncertainty flags its own failures.

An earlier workshop paper by the first author, in the MICCAI UNSURE 2026 proceedings, evaluated geometric TTA only on in-distribution data with frames split at random \cite{kashyap2026unsure}. This paper retrains every model under sequence- and patient-level splits, adds external validation, lesion-free populations and three ways of counting a miss, and carries over no figure, table or number; Appendix~\ref{app:further} compares the two directly, including where that study's conclusions do not replicate here.

\paragraph{Contributions.}
TTA with four or eight views raised Dice for every model on every dataset with lesions. Under shift, whether it lost lesions depended on the count, and the disagreement survives taking the checkpoint rather than a majority vote as the unit: UNet left more images empty in all 15 checkpoints on CVC-ColonDB and missed fewer at IoU 0.10 in all 15 on ETIS-Larib, and SwinUNETR's images with no overlap rose in all 15 on all three collections while its missed images moved on none. Requiring IoU 0.10 we could not detect a rise anywhere, but the widest interval still admits 20 more missed images in 1000, so that is a failure to detect harm and not a demonstration of safety. For empty outputs TTA was close to the curve traced by changing the threshold of a single view; what that change did not give was localisation. The fraction of pixels on which TTA views disagree did not flag empty predictions, scoring AUROC 0.21 and 0.03, below chance, and two rules reusing the same views filled most empty masks but fired on 94.5\% and 89.5\% of polyp-free frames.
